% concatenated from d2e-arxivv2.tex
%
%
%
\documentclass[runningheads, letter]{llncs}
\usepackage{hyperref}

%
%
\usepackage[ruled, lined, linesnumbered]{algorithm2e}
\SetKwProg{Try}{try}{:}{}
\SetKwProg{Catch}{catch}{:}{end}
\usepackage{algpseudocode}
\usepackage{graphicx}
\usepackage{extarrows}
\usepackage{amssymb,amscd}
\usepackage[all]{xy}
\usepackage{longtable}
%
\usepackage[latin1]{inputenc}
\usepackage{tikz,arydshln}
\usetikzlibrary{shapes,arrows}

\newcommand{\keywords}[1]{\par\addvspace\baselineskip
\noindent\keywordname\enspace\ignorespaces#1}

%
%
%
%
%
%
%
\newtheorem{conj}[theorem]{Conjecture}

\newcommand{\fq}{{\mathbb F}_{q}}
\newcommand{\ftwo}{{\mathbb F}_{2}}
\newcommand{\Fbar}{\overline{\mathbb F}}
\newcommand{\Gbar}{\overline{G}}
\newcommand{\diff}{\partial}
\newcommand{\prb}{\mbox{Prob}}
\newcommand{\pf}{{\em Proof.\ }}
\newcommand{\rmv}[1]{}

\newcommand{\N}{{\mathbb N}}
\newcommand{\Z}{{\mathbb Z}}
\newcommand{\R}{{\mathbb R}}
\newcommand{\F}{{\mathbb F}}
\newcommand{\C}{{\mathbb C}}
\newcommand{\Q}{{\mathbb Q}}
\renewcommand{\L}{{\mathcal L}}
\newcommand{\fp}{\mathfrak{p}}
\newcommand{\fP}{\mathfrak{P}}
\newcommand{\fa}{\mathfrak{a}}
\newcommand{\A}{{\mathbf A}}
\renewcommand{\P}{{\mathbf P}}
\renewcommand{\S}{{\mathbf S}}
\newcommand{\X}{{\mathbf X}}
\renewcommand{\b}{{\mathbf b}}
\newcommand{\B}{{\mathbf B}}
\renewcommand{\O}{{\mathcal O}}
\DeclareMathOperator{\Gal}{\mathrm{Gal}}

\pagestyle{plain}

%
\begin{document}
\mainmatter
\title
{Deterministic and Efficient Ideal Arithmetic via
  Two-Element Representations}
%
%

\author{Qi Cheng\inst{1} %
 }
 \institute{
   School of Computer Science\\
   University of Oklahoma\\
 Norman, OK 73019, USA.\\
 Email: {\tt qchengnorman@outlook.com}
\footnote{We thank the U. S. National Science Foundation for
  their support through the grant \texttt{CCF-2530361}.}
}
\maketitle

\begin{abstract}
  Given an ideal in a number field, it is desirable in many situations
  to find two elements that generate the ideal over the ring of the
  integers of the field.  Existing algorithms are either randomized,
  or impractical at cryptographic sizes.
  In the paper, we present a deterministic
  polynomial time algorithm to find the two-element representation of
  an ideal.  For a monic irreducible integral polynomial \( f(x) \),
  let \( K=\Q[x]/(f) \) be the number field, and
  \( O_K \) be the integral closure. Our algorithm works when
  the norm of the input ideal is co-prime to the index
  \( [O_K:\Z[x]/f] \). In particular, it
  handles all ideals for monogenic \( f(x) \),
  a class that includes the cyclotomic polynomials widely used in
  lattice based cryptography.  A key technical ingredient in our
  result is a generalized version of Dedekind criterion.
\end{abstract}

\section{Introduction}

Due to the imminent threat of quantum algorithms to
cryptosystems based on prime factorization and discrete logarithm
\cite{Shor94}, cryptographers are motivated
to design new cryptographic schemes that can resist quantum attacks.
Hard ideal lattice problems underpin the security of the Ring
Learning-With-Errors (RLWE) problem, a versatile and widely used
primitive in post quantum cryptography \cite{HPS98,LPR10,SSTX09}.
A large number of algorithms are proposed, including NewHope
\cite{NewHope}, Crystals-Kyber \cite{Kyber} and LAC \cite{LAC}.
%
%
These applications call for efficient implementation of
ideal arithmetic, an important topic in computational number
theory that was first studied in \cite{Cohen93,Cohen2012},
and the references therein.
%
%
It is well-known that for a number field,
two elements are sufficient to generate any fractional ideal as a
module over its ring of integers.  Several algorithms
for finding such representation have been proposed,
but they are either randomized or inefficient,
due to expensive prime factorization.
An ideal that admits a single generator is called principal. 
Interestingly there is no known efficient algorithm
to decide the principality of an ideal, let alone to find a
generator, even randomness is allowed.

\subsection{Description of the problem}

Let \( f(x) \) be a monic irreducible polynomial over \( \Z \)
of degree \( n \). Let \( K= \Q[x]/(f) \) be the number field.
We denote by \( O_K \)  the integral closure
of \( \Z[x]/f(x) \). It is a free \( \Z-\)module, but finding
its integral basis is in general difficult, often
requiring the prime factorization of the discriminant of \( f(x) \). 
An ideal in \( O_K \) is also a free \( \Z-\)module.  
If space complexity is not a primary concern,
then one \cite{Cohen2012}  can represent an ideal by its \( \Z-\)basis,
each of its elements being written in \( \Z-\)basis of \( O_K \).
This gives us an integral matrix (lattice) representation for an
ideal.

By sending \( x \) to one of the roots \( \alpha \)
of \( f(x) \) in \( K \), we establish an exact sequence
of \( \Z \)-modules: 
\[ 0 \rightarrow (f(x)) \rightarrow  \Z[x] \rightarrow O_K  \]
If the index \( [ O_K:\Z[x]/(f) ] \) is 1, then the field---
and the polynomial \( f(x) \)---
are called monogenic; the ring of integers has a power basis
\( \{1, \alpha, \alpha^2, \cdots, \alpha^{n-1}\} \);
and the sequence can be extended to \( 0 \):
\[ 0 \rightarrow (f(x)) \rightarrow  \Z[x] \rightarrow
  O_K \rightarrow 0. \]
This short exact sequence allows us to lift an ideal in \( O_K \) 
back to  an ideal in \( \Z[x]  \), and  compute in \( \Z[x] \),
which is often more efficient than  dealing with integral
matrix representations of ideals.
Moreover, ideals  \( I \subseteq \Z[x] \) and
\( I + f(x)\Z[x] \subseteq \Z[x]\)  
will be mapped to the same ideal in \( O_K \).
It means that we may freely add \( f(x) \) into an
ideal in \( \Z[x] \), without changing its image in \( O_K \).
Note that \( \Z[x] \) has ideals that can not be generated by
two elements (for instance \( (5, x)^3 \)).

In general, the problem of finding two-element representation
can be stated as follows: giving
\( a_1, a_2, \cdots, a_r  \) in \( O_K \),
to find \( b_1, b_2 \) in \( O_K \) so that
\[ a_1 O_K + a_2 O_K + \cdots + a_r O_K = b_1 O_K + b_2 O_K  \]
For example, if the ideal is given by its \( \Z \)-basis,
then \( r=n \).
We can assume that \( r=3 \), since \( a_4, a_5, \cdots \) 
can be absorbed into two elements recursively.
We can find an integer in the input ideal and assign it to
\( a_1 \). One of the two generators of the output ideal
can also be an integer.
%
So we should go one step further, and assume that
both \( a_1 \) and \( b_1\) are integers, and reformulate
the problem as:

\begin{quote}
 {\bf Input:} \( {\cal N} \in \Z, a(x) \in \Z[x], b(x)  \in \Z[x]\)

 {\bf Output:} \( M \in \Z, c (x) \in \Z[x] \) so that
 \[ {\cal N} O_K + a (\alpha) O_K + b (\alpha) O_K
   = M O_K + c(\alpha) O_K  \]
\end{quote}

In this paper, we focus on the most difficult case
where \( {\cal N} \) 
is a large composite number free of small prime factors.
Once we have an efficient algorithm for the above problem,
ideal addition is straight-forward:
\[ (N_1, a(x)) + (N_2, b(x)) = ( gcd(N_1, N_2), a(x), b(x) )  \]
Ideal multiplication and powering are core subroutines in
our algorithm, and they will be treated in
Lemma~\ref{lem:relativeprime} and
Section~\ref{sec:powerandmulti}.
%

\subsection{Previous works}

It has long been observed that two-element representations
can speed up computation significantly.
Pohst and Zassenhaus \cite{PZ85} utilized them
to calculate class groups.
Since their work dealt with ideals of small norms,
they can afford to compute two-element representations
by factoring ideals into  small prime ideals.
See also \cite{GMN13}.
It is difficult to generalize their method to ideals of
large norm, due to a few apparent difficulties here.
First the norm of the ideal need to be factored into primes,
which is hard. Second one must find  prime ideals
lying above a rational prime, which  requires to factor
polynomials over finite fields. It relies on  randomized algorithms
in the current state of the art \cite{KedlayaUm11}.

Another deterministic algorithm was sketched
in \cite[Subsection 6.3.2]{Belabas04}.
However, the time complexity was not analyzed, and
the authors deemed it impractical if
the ideal can not be factored.

%
%
A simple randomized algorithm was proposed in
\cite{MMPT23,PS21,FS10}.
The algorithm first finds \( a\in I \),
and then pick a random element \( r \in I/(a) \).
%
It can be proved that with high probability,
\( I = a O_K + r O_k \).
%


\subsection{Our result}

Assume that the input ideal is \( ({\cal N}, a(\alpha), b(\alpha)) \).
We address small prime factors of \( {\cal N} \) separately
via the Pohst-Zassenhaus algorithm.
Factoring polynomials in small finite fields does not
require randomness. So we may assume that \( {\cal N} \)
has no small prime factors (and is thus hard to factor).

{\bf Stage 1}
In the first stage of our algorithm, we factor the input ideal
into a product of ideals in two-element representation,
and ideals of a special form, described below.
The basic idea is to try to compute \( gcd(a(x), b(x)) \) 
over \( \Z/{\cal N} \) using the Euclidean algorithm. 
%
%
%
If the computation succeeds, we find a polynomial \( c (x) \),
and obtain two-element representation
\( {\cal N} O_K + c(\alpha) O_K \).
This should solve the problem for many inputs if \( {\cal N} \)
is free of small prime factors.
%
%

An obstacle in this step is a zero divisor in \( \Z/{\cal N} \)
appearing as the leading coefficient of some polynomials.
If additionally \( {\cal N} \) is not a perfect power,
we split \( {\cal N}  \) into factors which are co-prime
to each other.
%
This allows us to split the input
ideal. We then work on each factor ideal separately.
%

A degenerate case, when \( {\cal N} = N^e \) 
is a perfect power, must be handled separately.
For this case, we  run the Euclidean algorithm
on \( a(x) \) and \( b(x) \) over \( \Z/{N} \)
(rather than over \( \Z/{\cal N} \)), 
producing an ideal of form
\begin{equation}\label{eq:onegennotN}
 (N^e, a_1 (x), N d_1(x), N d_2(x), \cdots ),
\end{equation}
where only one  generator is not a multiple of
\( N \). We add \( f(x) \) into the ideal, and
it guarantees that \( a_1(x) \) is a factor of \( f(x) \)
over \( \Z/N \).

In the end, we factor the input ideal into
a product of factors in special forms:
\[ ({\cal N}, a(\alpha), b(\alpha))=
  \prod_{i=1}^t (N_i, a_i (x)) \prod_{i=t+1}^s (N_i^{e_i}, a_i (x),
  N_i d_{ 1 i}(x), N_i d_{ 2 i}(x), \cdots).\]

{\bf Stage 2}
The second stage of our algorithm will handle
ideals in the form of (\ref{eq:onegennotN}).
One can clearly read off a factor \( (N, a_1(\alpha)) \)
from the ideal that is in two-element representation.
The other factors will also likely contain \( N \).
%
How do we multiple \( (N, a_1(\alpha)) (N, a_2(\alpha)) \)
into a two-element form? Preferably the product is
\( (N^2, a_1(\alpha) a_2(\alpha)) \),
but this is not always true.
\begin{example}\label{ex:1}
  Consider the 8-th cyclotomic field with ring of
  integers \( \Z[x]/(x^4+1)  \).
  %
  The product of ideals \( (65, x^2 + 268) \) and \( (65, x^2 - 268) \)
  is:
  \begin{align*}
    & (65^2, 65(x^2+268), 65(x^2-268), (x^2+268)(x^2-268)) \\
    =& (65^2, 65 * 2 * 268, 65(x^2-268), -71825)\\
    =& (65),
  \end{align*}
  but
\[  (65^2, (x^2+268)(x^2-268)) = (65^2, -71825)
  = (65^2, -17 * 65^2) = (65^2). \]
  If we add  some
  multiplies of \( 65 \) to the second generator,
  \[ (65, x^2 + 268)=(65, x^2+8), (65, x^2 - 268)
    = (65, x^2-8),\]
  then
\[  (65^2, (x^2+8)(x^2-8)) = (65^2, -65) = (65). \]
%
%
The main innovation of our paper is to find 
the desirable second generator so that ideal arithmetic can
be done efficiently.
%
\end{example}
%
In the work of Pohst-Zassenhaus, when the integer generator
\( N \) is a prime, a special type of two-element representations
was found so that multiplication can be done
in a straight forward way:
\[ (p, \beta_1) (p, \beta_2) = (p^2, \beta_1 \beta_2 ). \]
%
%
Their algorithm is not available when \( N \) is not a prime,
since it requires prime ideal factorization.
Our main technical contribution is to  derive such a form
for a composite \( N \) with unknown prime factorization,
%
%
%
 overcoming the inefficiency in
 \cite[Subsection 6.3.2]{Belabas04}.
 See Section~\ref{sec:dedekind} and
\ref{sec:powerandmulti}.
%
%
%
%
%
%

{\bf Stage 3} After the second stage,
the algorithm splits the input ideal uniformly:
\[ ({\cal N}, a(x), b(x)) = \prod_{i=1}^s (N_i, a_i (x)), \]
where all the \( N_i \)'s are pairwise co-prime to each other.
Multiplying them together into a single ideal
in two-element representation
can be accomplished in a standard way by Chinese Remainder Theorem,
which is the goal in the final stage.
 See Lemma~\ref{lem:relativeprime}.
%
%
%

\begin{theorem}
  (Main) There is an algorithm solving the
  two-element representation problem for a monogenic field
  of degree \( n \) in time \( poly(n, \log {\cal N}) \).
  %
  %
  %
  When \( {\cal N} \) has no repeated prime factor, 
  the algorithm runs in time \( O(n^2 \log {\cal N} \omega({\cal N})) \),
  where \( \omega ({\cal N}) \) denotes the number of distinct prime
  factors of \( {\cal N} \).
\end{theorem}

Our algorithm shows that neither randomness nor prime
ideal factorization is necessary to find
two-element representation for an ideal.
In fact, finding a factor of \( {\cal N} \) only slows
down our algorithm.
Our algorithm is practical, since the basic operation is
arithmetic on polynomials rather than integral matrices
that is much more expensive.  Note that \( \omega({\cal N}) \)
is less than \(  \log {\cal N} \), and has average order
\( \log\log {\cal N} \)\cite{HW08}.
Thus when \( {\cal N} \) is square-free,
the complexity is comparable to 
randomized algorithm. Our algorithm
  works for all ideals in monogenic fields, covering
  cyclotomic fields---the most important case
  for ideal-based cryptography.
It also works for non-monogenic fields if the norm of I is
relatively prime to the index of $O_k$
over $\Z[x]/(f)$. Moreover, if the algorithm fails,
it will provide a certificate that \( f(x) \)
is not monogenic.
%
%
%
%


\section{Mathematical Preparation}

It is well-known \cite{Cohen93} that
a prime ideal containing \( p \) has two-element
representation \( (p, t(x)) \) if
\( p \) does not divide the index of \( [O_K:\Z[x]/(f)] \).
 Here \( t(x) \) is an irreducible factor of \( f(x) \)
over \( \F_p \).

In the section, we present a few facts about ideals
in number fields. They are known in different forms in literature.
We include some proofs here for completeness.

\subsection{The minimal positive integer in an ideal}

Given an ideal \( I \) in \( O_K \), we denote the
minimal positive integer in the ideal by \( MPI(I) \).
It divides the norm and can be computed by
Hermite Normal Form of the lattice corresponding to
\( I \). If the number field is monogenic and
the ideal is principal \( (a(x)) \), it can
be also computed by more efficient Euclidean algorithm. 
First we find the inverse of \( a(\alpha) \) in the field,
by running the Euclidean algorithm on \( f(x) \) and \( a(x) \)
in \( \Q[x] \). The inverse must be in the form
\( b(\alpha) \) for some polynomial in \( b(x) \in \Q[x] \).
Let the number \( i\) be the smallest positive integer
such that \( i b(x) \in \Z[x] \).
It can be shown that \( i =  MPI (a(\alpha))  \).
If an ideal is generated by a few elements \( \alpha_1,
\alpha_2, \cdots\), we can set \({\cal N} \) to be
\( gcd(MPI((\alpha_1)), MPI((\alpha_1)), \cdots) \).
%
%
%

%
%
%
%
%
%
%

\subsection{When integer generators are relatively prime}

We now show that factorization of an integer
generator of an ideal leads to factorization of the ideal.

\begin{lemma} \label{lem:factorideal}
  Let \( I \) be an ideal in \( \Z[x] \). Suppose that
  an integer \( {\cal N} \) factors as \( {\cal N} = N_1 N_2 \),
  with \( (N_1, N_2) = 1 \). Then the ideal factors as 
  \[ {\cal N} \Z[x] + I = (N_1 \Z[x]+ I) (N_2 \Z[x] + I)  \] 
\end{lemma}

\begin{proof}
  Since 
  \[  (N_1 \Z[x]+ I) (N_2 \Z[x] + I) = N_1 N_2\Z[x]
  + N_1 I \Z[x] + N_2 I \Z[x] + I^2, \]
it is clear that \(  RHS \subseteq LHS \).
   To see that
  \( LHS \subseteq RHS \), we observe
  that for any \( f(x) \in I \),
  \( f(x) \in f(x) N_1 \Z[x] + f(x) N_2 \Z[x] \). 
\end{proof}

\begin{lemma}
\label{lem:relativeprime}
  Assume that two positive integers \( N_1, N_2 \)
  are relatively prime. Let
  \[ M_1 =  N_1^{-1} \pmod{N_2 }\ and
  \ M_2 =  N_2^{-1} \pmod{N_1 }.\]
Then two ideals in \( \Z[x] \) can be multiplied
in polynomial time as
  \[ (N_1, f_1(x)) (N_2, f_2(x)) = (N_1 N_2, g(x) ), \]
  where \( f_1(x), f_2(x) \) are in \( \Z[x] \), and 
  \( g(x) = N_2 M_2 f_1(x) + N_1 M_1 f_2(x) \).
\end{lemma}

\begin{proof}
Essentially we find a polynomial \( g(x) \) so that
\( g(x) \pmod{N_1} = f_1(x) \)
and \( g(x) \pmod{N_2} = f_2(x) \), using the
Chinese Remainder Theorem.
Then we apply Lemma~\ref{lem:factorideal}.
\end{proof}


\section{An Euclidean algorithm for integral polynomials}

First we rephrase the Euclidean algorithm for polynomials
over \( \Z/N\Z \) as ideal normalization in \( \Z [x] \).

\begin{lemma}\label{lem:polygcd} 
  Let \( N \) be a positive integer, %
  and \( a(x), b(x) \) are two integral polynomials,
  neither of which is a multiple of \( N \).
  There is a polynomial time algorithm that either finds
  a nontrivial factor of \( N \), or two polynomials
  \( c(x) \) and \( d(x ) \) in \( \Z[x] \)  such that
  \[ a(x)\Z[x] + b(x)\Z[x] = c(x)\Z[x] + N d(x)\Z[x].  \]
  %
  In addition, in this case the algorithm also find three polynomials
  \( a'(x), b'(x), t(x) \)  such that
  \[ a(x) a'(x) + b(x) b'(x) + N t(x) = c(x).  \]
  %
\end{lemma}

\begin{proof}
  %
  %
  If \( N \) is a perfect power, we can find its non-trivial factor
  in polynomial time. Now assume that \( N \) is not a perfect power.

  If any coefficients of \( a(x), b(x) \) is a
  non-trivial zero divisor in \( \Z/N \), we can find
  a non-trivial factor of \( N \).
  Otherwise we separate the terms of \( a(x), b(x) \):
  \[ a(x)=a_1 (x) + N a_2(x), b(x) = b_1(x) + N b_2(x),  \]
  where the leading coefficients of \( a_1(x), b_1(x) \)
  are in \( ( \Z/N )^* \). W.l.o.g, assume that
  \( \deg(a_1(x)) \geq \deg (b_1(x)). \)
  Then in \( Z/N[x] \), we can divide \( a_1(x) \)
  by \( b_1(x) \), and let \( q(x) \in \Z[x] \) be the quotient.
  It implies that there exists a polynomial \( r(x) \in \Z[x] \)
  such that 
\[ a(x) = q(x) b(x) + r(x), \]
and \( r(x) = r_1(x) + N r_2(x)\),
\( deg(r_1 (x)) < deg(b_1 (x)) \).
As ideals in \( \Z[x] \), we have
\[ a(x) \Z[x] + b(x) \Z[x] = b(x) \Z[x] + r(x) \Z [x],  \]
since \( r(x) \) is in the LHS ideal
and \( a(x) \) is in the RHS ideal.

At this point, if any coefficient of \( r(x) \)
is a non-trivial zero divisor of \( \Z/N \),
then we find a non-trivial factor of \( N \).
Otherwise we divide \( b(x) \) by \( r(x) \),
and continue an Euclidean-style algorithm,
until the remainder polynomial is \( 0 \pmod{N} \).
This concludes the proof.
\end{proof}

The algorithm is basically the Euclidean algorithm in
\( \Z/N[x] \), which is implemented in many software
packages such as Sagemath. If the algorithm terminates without
encountering a non-trivial zero divisor in \( \Z/N \),
then we solve the two-element
representation problem for the ideal
\( (N, a(\alpha), b(\alpha)) \):
\[ N O_k + a(x) O_k + b(x) O_k = N O_K + c(x) O_k.  \]
What if a non-trivial zero divisor in \( \Z/N \) is found?
%
We also need to take care of a degenerate case,
when \( N \) is  a perfect power.
The following theorem shows that  we can
always factor and normalize the input ideal:

\begin{theorem}
  There is a polynomial time algorithm that,
  given two integral polynomials \( a(x ) \) and
  \( b(x ) \), a positive integer \( N  \)
  that is not a perfect power, and a positive integer \( e \),
  factors the ideal \( (N^e, a(x), b(x)) \subseteq \Z[x] \)
  into a product as
  \[ (N^e, a(x), b(x)) = \prod_{i} (N_i^{e_i}, a_i(x), N_i d_i (x))  \]
  such that none of \( N_i \)'s is a perfect power,
  and for any \( i \not= j \), \( gcd(N_i, N_j) = 1\).
\end{theorem}

\begin{proof}
  Run the algorithm from Lemma~\ref{lem:polygcd}
  on \( a(x), b(x) \) and \( N \) ( rather than \( N^e \) ).
  If we find \( c(x) \) and \( d(x) \) such that
  \( a(x)\Z[x] + b(x) \Z[x] = c(x)\Z[x] + N d(x) \Z[x], \)
  then
  \[  (N^e, a(x), b(x)) = (N^e, c(x), N d(x)). \]
  Otherwise we find a factorization of \( N \), then
  \[ (N^e, a(x), b(x)) = (N_1^{e_1}, a(x), b(x))
      (N_2^{e_2}, a(x), b(x)).\] 
    Here \( N_1, N_2 \)  are co-prime.
    We then recursively work on each of the factors,
    ultimately obtaining the desired factorization.
\end{proof}

We can add \( f(x) \) into the ideal,
%
run the
algorithm from Lemma~\ref{lem:polygcd} on \( a_i(x) \),
\( f(x) \) and \( N \),
and derive a corollary:

\begin{corollary}
\label{cor:gcd}
  There is a polynomial time algorithm that,
  given two integral elements \( a(\alpha ) \) and
  \( b(\alpha) \) in \( O_K \), a positive integer \( N  \)
  that is not a perfect power, and a positive integer \( e \),
  factors
  the ideal \( (N^e, a(\alpha), b(\alpha)) \subseteq O_K \)
  into a product as
  \begin{equation}
\label{eq:5}
 (N^e, a(\alpha), b(\alpha)) =
 \prod_{i} (N_i^{e_i}, a_{ i }(\alpha),
 %
    N_i I_i)  
  \end{equation}
  such that none of \( N_i \)'s is a perfect power;
  \( I_i = (d_{i1} (\alpha), d_{i2} (\alpha)) \)
  is an ideal in \( O_K \); and
  for any \( i \not= j \), \( gcd(N_i, N_j) = 1\);
  and \( a_{i}(x) \) is a polynomial dividing \( f(x) \)
  in \( \Z/N_i[x] \).
\end{corollary}

\begin{algorithm}[h]\label{alg:stage1}
  \caption{Sketch of the algorithm in the first stage}
    \KwIn{ \( f(x) \) monogenic, an integer \( {\cal N}\in \Z \)
    and two elements \( a(x), b(x) \) in \( \Z[x] \).}
  \KwOut{ Two stacks: R1 and R2.  }
  \tcc{Elements in R1 are pairs
    of form \( (N_i, a_i(\alpha)) \), elements in R2 have form
    \( ( N_i^{e_i}, a_i(\alpha),
    N_i d_{i1}(\alpha), N_i d_{i2}(\alpha),\cdots ) \),
    and all \( N_i \) are relatively co-prime.
    The product of all ideals in R1 and R2 is equal to
  \( ({\cal N}, a(\alpha), b(\alpha)) \)}
  Initiate an empty stack S\;
  Push \( {\cal N} \) into S\;
  Initiate an empty stack R1\;
  Initiate an empty stack R2\;
  \While{S is not empty}{
    \Try{}{
      Pop an integer  from S, and write it as \(N^e\) so
      that \( N \) is not a perfect power\; 
      \tcc{\( N^e \) is of course a perfect power if \( e \geq 2 \)}
      Run the algorithm from Lemma~\ref{lem:polygcd}
      on \( a(x), b(x)\) and \( N  \) to find
      \( c(x), d(x) \) so that over \( \Z[x] \),
      \( a(x) \Z[x] + b(x) \Z[x] = c(x)\Z[x] + N d(x)\Z[x]  \)\;
      \uIf{\( e = 1 \)}
        {Push \( (N^e, c(\alpha)) \) into R1\;
        }
      \Else{Run the algorithm  from Lemma~\ref{lem:polygcd}
          on \(c(x)\), \(f(x)\) and \( N \) to find
      \( a_1(x), d_2(x) \) so that over \( \Z[x] \),
      \( c(x) \Z[x] + f(x) \Z[x] = a_1 (x)\Z[x] + N d_2(x)\Z[x]  \)\;
          push \( (N^e, a_1 (\alpha), N d(\alpha), N d_2 (\alpha)) \) into R2; 
        }
    }
    \Catch{\( N \) is factored}
    {\For{each factor \( N_i \) with
          \( gcd(N, N/N_i) =1 \)}{
      Push \( N_i^{e}\) into the stack S\; }
    }
    }
\end{algorithm}

The first stage of the algorithm will factor the input
ideal as in Equation~\ref{eq:5}, and is described in pseudo-code
in Algorithm~\ref{alg:stage1}. 
Note that we may write \( a_{i} (x) \), in many ways,
as a monic polynomial dividing \( f(x) \) over \( \Z/N_i \), plus
a multiple of \( N_i \) in \( \Z[x] \).
%
%
%
%
This flexibility plays an essential
role in later stages of the algorithm.
%
%
%
%
For details see Theorem~\ref{thm:tech}. 
For a factor in the RHS of Equation~\ref{eq:5} with \( e_i =1 \),
it is already in the two-element form.
While other  factors in the RHS with \( e_i > 1 \)
is not in two-element representation,
we can immediately read off a factor \( (N_i, a_{i}(\alpha)) \)
with two generators.
%
%
If \( a_{i}(\alpha) = 1\pmod{N} \), then
\(  (N_i^{e_i}, a_{ i }(\alpha), N_i I_i)\) is
    the entire ring \( O_K \), according to the following lemma.

\begin{lemma}
  If \( A \in ( \Z/N )^* \), then for any \( g(x) \),
  \( A+N g(x) \) is a unit in \( \Z/N^s [x] \) for
  any positive integer \( s \).
\end{lemma}

\begin{proof}
  Let \( B \in \Z \) be an inverse of \( A \), namely,
  there exists an \( C\in \Z \) such that \( A B = 1 + N C \).
  We have
  \begin{align*}
 1/(A + N g(x)) =& B/ ( A B + N B g(x))\\
    =& B/ ( 1 + N (B g(x) + C) )\\
    =& B ( 1 - N (B g(x) + C) + (N (B g(x) + C))^2 \\
    & \cdots + (- N (B g(x) +C))^{s-1} ) \pmod{N^s }
  \end{align*}
\end{proof}

If \( a_{i}(\alpha) \not= 1 \pmod{N_i} \),
we need to compute
\( (N_i^{e_i}, a_{ i }(\alpha), N_i I_i)/(N_i, a_{i}(\alpha)) \).
%
In the next two sections, we show how to
find a proper member in  \( a_{i} (x) + N_i \Z[x] \) 
so that ideal division, multiplication and powering
behave nicely.

\section{Generalized Dedekind's criterion}
%
\label{sec:dedekind}

Let \( f(x) \) be monogenic.
Let \( a_1(x ) \in \Z[x] \) be a monic factor
of \( f(x) \) over \( \Z/N \), where \( N \) is not a perfect power.
We can write
\begin{equation}\label{eq:1} 
 f(x) = f_1 (x) N^s + a_1^e (x) A(x),  
\end{equation}
such that \( f_1(x) \not= 0 \pmod{N} \), 
and \( a_1(x) \) does not divide \( A(x)\) (which is also monic)
over \(\Z/N \).

\begin{lemma}
\label{lem:gended}
If \( e\geq 2 \), we must have \( s =1 \).
Moreover, if running the algorithm from Lemma~\ref{lem:polygcd}
on \( f_1(x), a_1(x) \)
and \( N \) does not throw an exception of
finding a non-trivial factor of \( N \),
  then \( (f_1(x), a_1(x), N) =1 \) in \( \Z[x] \).
\end{lemma}

\begin{proof}
  If  \( s \geq 2 \), or if \( a_1 (x) \)  shares
  factors with \(  f_1(x) \) over \( \Z/N \),
  then it violates the Dedekind's criterion
\cite[page 31]{GW21}\cite{Dedekind1878}
  at \( (p, a'(x)) \),
  where \( p \) is a prime factor of \( N  \),
  and \( a'(x) \) is an irreducible factor of \( a_1(x) \pmod{p} \). 
  In other words,
  \( (p, a'(x)) \) is a maximal ideal corresponding
  to a singular point on the
  arithmetic curve \( f(x) \) over \( \Z \).  
\end{proof}

The above lemma is basically a generalized version of Dedekind's criterion.
We will use it to prove the core technical lemma of this paper.

\begin{lemma}
  %
  %
  Assume that \( e=1 \).
  In polynomial time, we can either find a non-trivial
  factor of \( N \),
  or find two monic integral polynomials \( a_2 (x) \)
  and \( a_3 (x) \),
  such that \( f(x) \) can be rewritten as  
  \begin{equation} \label{eq:2}
   f(x) =f_2(x) N^{s_2} + a_2 (x) a_3 (x) A(x); 
  \end{equation}
  \( a_1 (x)=a_2 (x) a_3 (x) \pmod{N} \);
  \( (a_3 (x), f_2(x))_N =1 \); \( a_2(x) \)
  divides \( f_2(x) \) over \( \Z/N \);
  and \( (a_2(x), A(x), N) = 1\) in \( \Z[x] \).
\end{lemma}

\begin{proof}
  %
  %
  %
  Starting from Equation~(\ref{eq:1}),
  we run the algorithm from Lemma~\ref{lem:polygcd}
  on \linebreak \( (a_1(x), f_1 (x)) \) and \( N \). 
  If \( a_1(x) \) divides \( g(x) \),
  or the gcd is 1, then we are done. 
  Otherwise, we obtain a factorization of \( a_1(x) \) 
  over \( \Z/N \), and rewrite \( f(x) \) as
  \[ f(x) = f_3(x) N^{s_3} + a_4 (x) a_5 (x) A(x). \]
  Note that \( f_3(x) \) may differ from \( f_1(x) \).
  We run the
  algorithm from Lemma~\ref{lem:polygcd} on \( a_4(x), f_3(x), N \) and
  \( a_5(x), f_3(x), N \), and continue to find factors of \( a(x) \)
  over \( \Z/N \).
  %
  %
  We terminate the process when
  \[ f(x) = f_2 (x) N^{s_2} + ( \prod_i a_{i2}(x)) (\prod_i a_{i3} (x))
    A(x), \]
  each of \( a_{i2}(x) \)  divides \( f_2 (x) \) over \( \Z/N \),  
  and each of \( a_{i3}(x) \) is co-prime to
  \( f_2 (x) \) over \( \Z/N \).  Since \( a_1 (x) \) can be split
  into at most \( \deg(a_1(x)) \) factors, the
  algorithm will terminate eventually.
  The lemma follows by setting
  \[ a_2(x) = \prod a_{i2} (x), a_3(x) = \prod a_{i3}(x).  \]
  Note that \( a_2(x)   \) can not share a factor with \( A(x) \) 
  by the generalized
  Dedekind's criterion. This is crucial for the proof.
\end{proof}

%
%
Although we have found some factors of
\( a_1 (x) \) over \( \Z/N \),  
there is in general no efficient way to factor 
polynomials when \( N \) is not a prime.
In fact, even taking square roots modulo \( N \) is
as hard as factoring \( N \).

\begin{theorem} \label{thm:tech}
  Let \( a(x) \) be a factor of \( f(x) \) over \( \Z/N \).
  %
  In polynomial time, we can either find a non-trivial
  factor of \( N \),
  or find a monic integral polynomials \( \tilde{a} (x) \)
  such that \( f(x) \) can be written as  
  \begin{equation}
\label{eq:3}
 f(x) =\tilde{f_1}(x) N + \tilde{a}(x) \tilde{A}(x), 
  \end{equation}
  so that \( a(x)=\tilde{a} (x)  \pmod{N} \),
  and \( (\tilde{a} (x), \tilde{f_1}(x), N)=1 \) in \( \Z[x] \).
\end{theorem}

\begin{proof}
  The case of \( e\geq 2 \) is covered by Lemma~\ref{lem:gended}.
  Now assume that \( e=1 \).
  We first find \( a_2(x) \) and \( a_3(x) \) in Equation~(\ref{eq:2}).
  If \( s\geq 2 \), then by Lemma~\ref{lem:gended},
  \( a_2(x) a_3(x) \) can not share a factor with
  \( A(x) \), so we write
  \[ f(x) = (f_2(x) N^{s-1} -   A(x)) N +
    ( a_2 (x)  a_3 (x)+ N ) A(x), \]
  and set \( \tilde{a}(x) =  a_2 (x)  a_3 (x)+ N \)
  and \( \tilde{f_1}(x) = f_2(x) N^{s-1} -   A(x) \).
  If \( s= 1 \), we have
  \[ f(x) = (f_2(x)  -  a_3(x) A(x)) N +
    ( a_2 (x) + N ) a_3 (x) A(x), \]
  and set \( \tilde{a}(x) = ( a_2 (x) + N ) a_3 (x) \)
  and \( \tilde{f_1}(x) = f_2(x)  -  a_3(x) A(x) \).
  In either case
  it can be verified that the resulting form satisfies
  the requirements in the theorem.
\end{proof}

\section{Powering/Multiplying ideals in two-element representations}
\label{sec:powerandmulti}
We now return to an ideal
factor in the RHS of Equation~(\ref{eq:5}).
We will omit the subscription for convenience,
and denote it by \( J =  (N^e, a(\alpha), N I ) \). 
Recall that \( N \) is not a perfect power.
We find \( \tilde{a}(x) \) satisfying
the condition in Equation~(\ref{eq:3}),
and read off a factor \( (N, \tilde{a} (\alpha)) \) from \( J \).
We also find \( g(\alpha) \)
so that \( a(\alpha) = \tilde{a}(\alpha) + N g(\alpha) \).
%
%
According to Theorem~\ref{thm:tech}, we may write
\begin{equation}\label{eq:4} 
 f(x) = f_1 (x) N + \tilde{a} (x) A(x).  
\end{equation}
with \( (f_1(x), \tilde{a}(x), N) = 1 \) in \( \Z[x] \)
(which implies that \( f_1(x) \not= 0 \pmod{N} \) ). 
We will show that the condition will enable fast ideal arithmetic.

\begin{example}
  In Example~\ref{ex:1}, we can see that it is
  easier to multiply when using \( a(x) = x^2+8 \),
  rather than \(x^2 + 268\), as the second generator,
  even though \( (65, x^2 + 8) = (65, x^2 + 268) \).
  It is because that
  \[ f(x)= x^4+1 = (x^2+8) (x^2-8)+65  \]
  so \( x^2+8 \) satisfies the conditions in
  Equation~(\ref{eq:4}), whereas \( x^2+268 \) does not,
  as
  \[ f(x) = x^4+1 = (x^2+268) (x^2-268)+65^2*17.  \]
\end{example}

\begin{lemma}\label{lem:genN}
  For any positive integer \( s \), \( N \in (N^s, \tilde{a}(\alpha) )\).
  %
\end{lemma}

\begin{proof}
Since \( (\tilde{a}(x), f_1 (x), N ) = 1 \),
there must exist integer polynomials \( a'(x), f_1'(x) \)
and \( t(x) \) such that
\[ \tilde{a}(x) a'(x) + f_1(x) f_1'(x) + N t(x) = 1.  \]
We also have
\[ f_1 (\alpha) N + \tilde{a}(\alpha) A(\alpha) =0 \]
so
\[ f_1'(\alpha) f_1 (\alpha) N + f_1'(\alpha) \tilde{a}(\alpha) A(\alpha) =0 \]
\[ (1 - \tilde{a}(\alpha) a'(\alpha) - N t(\alpha)) N + f_1'(\alpha)
  a(\alpha) A(\alpha) =0 \]
\[ (1 -  N t(\alpha)) N = \tilde{a}(\alpha) a'(\alpha) N
  - f_1'(\alpha) \tilde{a} (\alpha) A(\alpha)  \]
\[ (1 -  N t(\alpha)) N = \tilde{a}(\alpha) ( a'(\alpha) N
  - f_1'(\alpha)   A(\alpha)  ) \]
The statement follows from the fact that
\[ 1- N t(\alpha)  \] is a unit in
\( \Z[\alpha]/(N^{s}) \).
\end{proof}

By setting \( \bar{a}(\alpha) = (a'(\alpha) N
- f_1'(\alpha)  A(\alpha))/(1 - N t(\alpha)) \in O_K \),
we conclude:

\begin{corollary}
  For \( \tilde{a} (x) \) in Equation~\ref{eq:4},
  we can find \( \bar{a}(\alpha) \) in deterministic
  polynomial time such that
  \[ \tilde{a}(\alpha)  \bar{a}(\alpha)  = N \pmod{N^s}  \]
  for any \( s\geq 2 \). Moreover
  \[ (N, \tilde{a}(\alpha)) (N, \bar{a}(\alpha)) = N   \]
\end{corollary}

Now we are ready to calculate \( J/ (N, \tilde{a}(\alpha))  \)
\begin{lemma}
\label{lem:other}
Let \( \tilde{a} (x) \) satisfy the conditions
in Equation~\ref{eq:4}.
Given an ideal \( J = (N^e, \tilde{a}(\alpha) + N g(\alpha), N I) \),
where \( I = (d_1 (\alpha), d_2(\alpha), \cdots, d_k (\alpha)) \),
there is a polynomial time algorithm that either finds a non-trivial
factor of \( N \), or calculate \( J/ (N, \tilde{a}(\alpha)) \),
and normalize it into a form
\( (N^e, a' (\alpha) , N I'), \)
where \( a' (x) \) is a factor of \( f(x) \)
over \( \Z/N \), and \( I' \) has at most \( k+2 \) generators.
\end{lemma}

\begin{proof}
  We have
\begin{align*}
  %
  &J/ (N, \tilde{a}(\alpha))\\
  =&  (N, \bar{a}(\alpha))(N^e, \tilde{a}(\alpha) +N g(\alpha), N I)/N\\
  =&  (N^{e+1}, N \tilde{a}(\alpha) +N^2 g(\alpha),
     N^2 I, N^{e} \bar{a}(\alpha),
      \tilde{a}(\alpha) \bar{a}(\alpha)+ N g(\alpha)\bar{a}(\alpha),
     \bar{a}(\alpha) N I)/N\\
  =&  (N^{e},  \tilde{a}(\alpha) +N g(\alpha),
     N I, N^{e-1} \bar{a}(\alpha),
     ( \tilde{a}(\alpha) \bar{a}(\alpha) )/N
     + g(\alpha)\bar{a}(\alpha),
     \bar{a}(\alpha)  I)\\
  =&  (N^{e},  \tilde{a}(\alpha) +N g(\alpha),
     ( \tilde{a}(\alpha) \bar{a}(\alpha) )/N
     + g(\alpha)\bar{a}(\alpha),
     \bar{a}(\alpha)  I, N^{e-1} \bar{a}(\alpha))
\end{align*}
Now we run the algorithm from Lemma~\ref{lem:polygcd} on
the polynomial generators that are not multiples of
\( N \), and normalize the above ideal into a form
\( (N^e, a'(\alpha) , N I') \)
where \( I' \) has at most \( k+2 \) generators.
\end{proof}

%
%
%
%
%
We could, of course, reduce the number of generators in \( I' \),
but it may not be worth the computation effort.
We should continue producing ideal factors and put them in
the preferred form.
%
%
The following theorem provides a simple algorithm
to multiply them into one ideal in two-element representation.

\begin{theorem}
  If \( \tilde{a}_1(x), \tilde{a}_2(x), \cdots, \tilde{a}_k(x) \) all
  satisfy the conditions in Equation~(\ref{eq:4}),
  then
  \[ (N, \tilde{a}_1(\alpha)) (N, \tilde{a}_2 (\alpha)) \cdots (N, \tilde{a}_k(\alpha))
  = (N^k, \tilde{a}_1(\alpha)\tilde{a}_2(\alpha)\cdots \tilde{a}_k(\alpha))\]
\end{theorem}

\begin{proof}
  Let \( 1\leq s < k \) be an integer,
  and \( i_1, i_2, \cdots, i_s, i_{s+1}, \cdots i_k \)
  be a permutation of \( 1, 2, 3, \cdots, k \).
  We prove the theorem by showing that 
  \( N^{k-s} \tilde{a}_{i_1}(\alpha) \tilde{a}_{i_2}(\alpha) \cdots \tilde{a}_{i_s}(\alpha)\), which is a term in the expansion of LHS,
  is in the RHS ideal. 
  By applying Lemma~\ref{lem:genN}, we observe that
  for \( s+1 \leq j \leq k \),
  \[ N \in  (N^k, \tilde{a}_{i_{ j }}(\alpha) ), \]
  and derive:
  \begin{align*}
     & N^{k-s} \tilde{a}_{i_1}(\alpha) \tilde{a}_{i_2}(\alpha) \cdots \tilde{a}_{i_s}(\alpha)\\
    \in& ( \tilde{a}_{i_1}(\alpha) \tilde{a}_{i_2}(\alpha) \cdots \tilde{a}_{i_s}(\alpha)  )(N^k, \tilde{a}_{i_{s+1}}(\alpha) )
          \cdots (N^k, \tilde{a}_{i_k}(\alpha) )\\
  \subseteq& (N^k, \tilde{a}_1 (\alpha)\tilde{a}_2(\alpha)\cdots \tilde{a}_k(\alpha))
  \end{align*}
\end{proof}

Note that we do not require \( \tilde{a}_i(x) \) to be
relatively prime to each other over \( \Z/N \).
This theorem provides mathematical foundation for stage 2
of our algorithm. For pseudo-code, see Algorithm~\ref{alg:stage2}.

\begin{corollary}
  For any positive integer \( s \), we have
  \[ (N, \tilde{a}(\alpha))^s = (N^s, \tilde{a}(\alpha)^s) \] 
\end{corollary}


\begin{algorithm}[h]\label{alg:stage2}
  \caption{Sketch of the algorithm in the second stage}
    \KwIn{ The stack R1 and  R2 from the first stage.}
    \KwOut{ Updated R1}
  %
    \While{R2 is not empty}{
    \Try{}{
      Pop an ideal \( J= ({N^{e}}, a(\alpha), {N} *I) \) from R2\; 
       \( Int \gets 1 \)\;
       \( Poly \gets 1 \)\;
       %
    \While{\( {a(\alpha)}\not= 1 \pmod{{N}} \)}
      {
       %
       %
       %
        Find \( \tilde{a}(x) \) for \( a(x) \) as in Theorem~\ref{thm:tech}\;
        \( Int \gets Int * \mathtt{N} \)\;
        \( Poly \gets Poly * \tilde{a}(\alpha) \)\;
        \( J \gets J/ (\mathtt{N}, \tilde{a}(\alpha)) \) 
         (Lemma~\ref{lem:other})\;
         Normalize \( J \) and update
         \( {a,  I} \) so that
         \( J = ({N^{e}}, {a(\alpha)} , {N} * {I}) \)\; 
    }
    Push the ideal \( (Int, Poly) \) into R1\;
    }
    \Catch{\( {N} \) is factored}
    {\For{each factor \( N_1^{e_1} \) of \( N \) with
          \( gcd(N_1, {N}^{\mathtt{e}}/N_1^{e_1}) =1 \)}{
          Push \( (N_1^{e e_1}, a(\alpha),
          {N_1}*  (N_1^{e_1 - 1} I )) \)
          into the stack R2\;
    \tcc{Since \(N_1 | \mathtt{N}\), the ideal is in the right form.}
     }
    }
    }
\end{algorithm}

\section{Put the algorithm together}

The third stage of our algorithm, described in
Algorithm~\ref{alg:stage3}, is a straight-forward
application of Lemma~\ref{lem:relativeprime}.
Given an input ideal, we run Algorithm~\ref{alg:stage1},
Algorithm~\ref{alg:stage2} and Algorithm~\ref{alg:stage3}
in sequence, obtaining two elements that generate
the input.
%
%
The building block of our algorithm is the Euclidean
algorithm for integral polynomials over \( \Z/N \),
where \( N \) is a factor of \( {\cal N} \) and
is not a perfect power. To analyze the complexity, we count
the number of gcd calculations.
Since small prime factors of \( {\cal N} \) have been
taken care efficiently by the Pohst-Zassenhaus algorithm,
we focus on \( {\cal N} \) free of small prime factors.
The most time-consuming part of the algorithm involves
handling perfect powers.
%
In the generic case when \( {\cal N} \) is free of
repeated factors, the algorithm will need
at most \( \omega({\cal N}) \) gcd calculations
---which is on average \( \log\log {\cal N} \)---
resulting in a complexity of
\( O (n^2 \log {\cal N} \omega({\cal N}) )\).
Even when \( {\cal N} \) has repeated factors,
the algorithm may not discover them as an exception
and may thus achieve similar complexity.

\begin{algorithm}[h]\label{alg:stage3}
  \caption{Sketch of the algorithm in the third stage}
  \KwIn{A stack R1}
      \KwOut{ two elements $M \in \Z$ and $c(x) \in \Z[x]$ }
    \While{R1 is not empty}{
      pop the stack \( n1, b1 \)
      \uIf{R1 is empty}
        {\Return \( (n1, b1) \)\;}
      \Else{
        Pop the stack to \( (n2, b2) \)\;
        \( n = n1 * n2 \)\;
        $m1 = n1^{-1} \pmod{n2}$\;
        $m2 = n2^{-1} \pmod{n1}$\;
        \( b = n2 * m2 * b1 + n1 * m1 * b2  \)\;
        Push \( (n,b) \) into R1\;
        }
   }
 \end{algorithm}

 \section{Conclusion and open problems}
In this paper, we present an efficient implementation of
ideal arithmetic for monogenic polynomials.
Cyclotomic polynomials are monogenic,
making our algorithms directly applicable to
the most important case for cryptography based on ideal lattices.
Our algorithm may be viewed as a
generalization of the Pohst-Zassenhaus algorithm
to composite integers.

For a polynomial \( f \) that is not
monogenic, the field \( \Q[x]/f \) may still be
monogenic. An open problem is to find such a monogenic
polynomial for the field in polynomial time.
Due to the existence of inessential
discriminant divisors, we may not be able to find a monogenic
polynomial for every field, however, a question can
certainly be raised about finding a polynomial \( f' \)
efficiently to  minimize the index \( [O_K:\Z[x]/f'] \).
Because inessential discriminant divisors are in general very small, resolving this question could lead to substantial gains in computational efficiency.

\begin{thebibliography}{10}
\providecommand{\url}[1]{\texttt{#1}}
\providecommand{\urlprefix}{URL }
\providecommand{\doi}[1]{https://doi.org/#1}

\bibitem{NewHope}
Alkim, E., Ducas, L., P{\"{o}}ppelmann, T., Schwabe, P.: Newhope without
  reconciliation. {IACR} Cryptology ePrint Archive  \textbf{2016}, ~1157
  (2016), \url{http://eprint.iacr.org/2016/1157}

\bibitem{Belabas04}
Belabas, K.: Topics in computational algebraic number theory. Journal de
  theorie des nombres de Bordeaux  \textbf{16}(1),  19--63 (2004).
  \doi{10.5802/jtnb.433}

\bibitem{Kyber}
Bos, J.W., Ducas, L., Kiltz, E., Lepoint, T., Lyubashevsky, V., Schanck, J.M.,
  Schwabe, P., Seiler, G., Stehl{\'{e}}, D.: {CRYSTALS} - kyber: {A}
  {CCA}-secure module-lattice-based {KEM}. In: 2018 {IEEE} European Symposium
  on Security and Privacy, EuroS{\&}P 2018, London, United Kingdom, April
  24-26, 2018. pp. 353--367 (2018). \doi{10.1109/EuroSP.2018.00032},
  \url{https://doi.org/10.1109/EuroSP.2018.00032}

\bibitem{Cohen93}
Cohen, H.: A Course in Computational Algebraic Number Theory. Springer-Verlag
  (1993)

\bibitem{Cohen2012}
Cohen, H.: Advanced Topics in Computational Number Theory. Graduate Texts in
  Mathematics, Springer New York (2012)

\bibitem{Dedekind1878}
Dedekind, R.: Ueber den zusammenhang zwischen der theorie der ideale und der
  theorie der hoheren congruenzen. Abhandlungen der Koniglichen Gesellschaft
  der Wissenschaften in Gottingen  \textbf{23},  3--38 (1878),
  \url{http://eudml.org/doc/135827}

\bibitem{FS10}
Fieker, C., Stehl{\'e}, D.: Short bases of lattices over number fields. In:
  Algorithmic Number Theory. pp. 157--173. Springer Berlin Heidelberg, Berlin,
  Heidelberg (2010)

\bibitem{GW21}
Gouvea, F.Q., Webster, J.: Dedekind on higher congruences and index divisors,
  1871 and 1878 (2021), \url{https://arxiv.org/abs/2107.08905}

\bibitem{GMN13}
Guardia, J., Montes, J., Nart, E.: A new computational approach to ideal theory
  in number fields. Foundations of Computational Mathematics  \textbf{13}(5),
  729--762 (2013)

\bibitem{HW08}
Hardy, G.H., Wright, E.M.: An Introduction to the Theory of Numbers. Oxford,
  sixth edn. (2008)

\bibitem{HPS98}
Hoffstein, J., Pipher, J., Silverman, J.H.: {NTRU:} {A} ring-based public key
  cryptosystem. In: Algorithmic Number Theory, Third International Symposium,
  ANTS-III. pp. 267--288 (1998)

\bibitem{KedlayaUm11}
Kedlaya, K.S., Umans, C.: Fast polynomial factorization and modular
  composition. SIAM J. Comput.  \textbf{40}(6),  1767--1802 (2011)

\bibitem{LAC}
Lu, X., Liu, Y., Zhang, Z., Jia, D., Xue, H., He, J., Li, B.: {LAC:} practical
  {Ring}-{LWE} based public-key encryption with byte-level modulus. {IACR}
  Cryptology ePrint Archive  \textbf{2018}, ~1009 (2018),
  \url{https://eprint.iacr.org/2018/1009}

\bibitem{LPR10}
Lyubashevsky, V., Peikert, C., Regev, O.: On ideal lattices and learning with
  errors over rings. In: Advances in Cryptology - {EUROCRYPT}. Lecture Notes in
  Computer Science, vol.~6110, pp. 1--23. Springer (2010)

\bibitem{MMPT23}
Micheli, G.D., Micciancio, D., Pellet-Mary, A., Tran, N.: Reductions from
  module lattices to free module lattices, and application to dequantizing
  module-{LLL}. In: Proc. of the Annual International Cryptology Conference
  (CRYPTO). Springer International Publishing (2023),
  \url{https://eprint.iacr.org/2023/886}

\bibitem{PS21}
Pellet-Mary, A., Stehl{\'e}, D.: On the hardness of the {NTRU} problem. In:
  Advances in Cryptology -- ASIACRYPT 2021. pp. 3--35. Springer International
  Publishing, Cham (2021)

\bibitem{PZ85}
Pohst, M., Zassenhaus, H.: Uber die berechnung von klassenzahlen und
  klassengruppen algebraischer zahlkorper. Journal fur die reine und angewandte
  Mathematik  \textbf{1985}(361),  50--72 (1985).
  \doi{doi:10.1515/crll.1985.361.50},
  \url{https://doi.org/10.1515/crll.1985.361.50}

\bibitem{Shor94}
Shor, P.W.: Algorithms for quantum computation: Discrete logarithms and
  factoring. In: 35th Annual Symposium on Foundations of Computer Science -
  FOCS. pp. 124--134 (1994)

\bibitem{SSTX09}
Stehl{\'{e}}, D., Steinfeld, R., Tanaka, K., Xagawa, K.: Efficient public key
  encryption based on ideal lattices. In: Advances in Cryptology - {ASIACRYPT}
  2009, 15th International Conference on the Theory and Application of
  Cryptology and Information Security, Tokyo, Japan, December 6-10, 2009.
  Proceedings. pp. 617--635 (2009). \doi{10.1007/978-3-642-10366-7\_36},
  \url{https://doi.org/10.1007/978-3-642-10366-7\_36}

\end{thebibliography}

\end{document}
